\documentclass[12pt]{article}  
\usepackage{amsmath}
\usepackage{amssymb}

\usepackage{booktabs}
\usepackage{array}

\usepackage{hyperref}
\providecommand{\href}[1]{\url{#1}}
\usepackage{iftex}
\usepackage[utf8]{inputenc}  
\usepackage{graphicx}
\setlength{\parindent}{0pt}
\setlength{\parskip}{6pt plus 2pt minus 1pt}
\usepackage[normalem]{ulem}
\usepackage[
]{geometry}

\usepackage{amsthm}
\usepackage{float}

\theoremstyle{plain}
\newtheorem{theorem}{Theorem}[section]
\newtheorem{lemma}{Lemma}
\newtheorem{corollary}{Corollary}
\newtheorem{proposition}{Proposition}
\newtheorem{conjecture}{Conjecture}

\theoremstyle{remark}
\newtheorem{remark}{Remark}
\newtheorem{note}{Note}
\newtheorem{claim}{Claim}

\theoremstyle{definition}
\newtheorem{definition}{Definition}[section]  
\newtheorem{condition}{Condition}
\newtheorem{problem}{Problem}[section]
\newtheorem{example}{Example}[section]

\renewcommand\qedsymbol{$\blacksquare$}  

\usepackage{mathtools}
\usepackage{listings}

\usepackage{framed}

\ProvidesPackage{iftex}

\ifxetex

\usepackage{unicode-math}
% \setmainfont{STIX Two Text}
% \setmathfont{STIX Two Math}
% \setmathfont[range={\mathcal,\mathbfcal},StylisticSet=1]{STIX Two Math}

% \setmonofont{JetBrains Mono}

\fi

\providecommand{\tightlist}{\setlength{\itemsep}{0pt}\setlength{\parskip}{0pt}} 

\usepackage{csquotes}
\title{Sequences}
\author{https://github.com/icprplshelp/}
\date{\today}

\begin{document}


\maketitle

{
\setcounter{tocdepth}{3}
\tableofcontents
}

\section{The basics}

A note about the epsilon characterization of the inf or the sup, it can
only really be used if we know for a fact that we are dealing with
something that is confirmed to be the inf or sup. Don't try to prove
that something is the sup / inf using the epsilon characterization.

A sequence \(S = \left\{ a_{n} \right\}_{n = 1}^{\infty}\) is like a
regular function \(f(n) = a_{n}\), except you can only ask its limit as
\(n \rightarrow \infty\). Sequences are discrete, so you can't take the
limit anywhere but \(\infty\).

Properties:

\begin{definition}[Bounded above]
A sequence \(S\) is bounded above
if every number in the sequence is less than some big number \(M\). For
example:

\[a_{n} = \frac{1}{n}\]

is bounded above.
\end{definition}

If a sequence \(S\) is bounded above, there exists a quantity that is
very similar to the maximum of \(S\), It is called the \textbf{supremum}
of \(S\), denoted \(\sup S\). A supremum is a value if you take an
\(\varepsilon\)-sized step to the left, you will bump into something
inside \(S\).

If a sequence \(S\) is bounded below, there exists a quantity similar to
the minimum of \(S\), called the infimum of \(S\), denoted \(\inf S\).
An infimum is a value in which if you add \(\varepsilon\), then you're
in \(S\).

For example:

\[S_{1} = \left\{ x:|x - 1| < 2 \right\}\]

Then \(\sup S_{1} = 3,\inf S_{1} = - 1\). This is like a 2-radius circle
centered at 1.

\[S_{2} = \left\{ 0.9,\ 0.99,\ 0.999,\ \ldots \right\}\]

Then \(\sup S_{2} = 1,\inf S_{2} = 0.9\).

\[S_{3} = \left\{ ( - 1)^{n}\left( 1 - \frac{1}{n} \right) \right\}\]

Then \(\sup S_{3} = 1\) and \(\inf S_{3} = - 1\).


\subsection{Limits for sequences to
infinity}

Same as regular functions. The definition is:

\[\forall\varepsilon > 0,\ \exists N,\ n > N \Rightarrow \left| a_{n} - L \right| < \varepsilon\]

If the limit exists, then \(S\) gets close to \(L\), so it is not going
to infinity, so it is bounded above and below. However, being bounded
above and below is not enough for the limit to exist -- which occurs if
the sequence is oscillating.


\subsection{Convergent, divergent}

If the limit of \(S = \left\{ a_{n} \right\}\) DNE, we say it diverges.
It could:

\begin{itemize}
\item
  Approach \(\infty\)
\item
  Oscillate
\end{itemize}


\subsection{Bounded above and increasing implies
convergence}

\begin{itemize}
\item
  Increasing means \(\forall n,\ a_{n + 1} > a_{n}\) (successive terms
  are larger)
\item
  Non-decreasing means \(\forall n,\ a_{n + 1} \geq a_{n}\) (nothing
  goes down as \(n\) goes up)
\item
  If \(S\) is increasing, then it's non-decreasing.
\end{itemize}

If \(S\) is bounded above and it's \textbf{non-decreasing,} then it
converges to its \textbf{supremum.} It will be stuck at the supremum. It
will be bounded by the least upper bound.

\begin{proof} Suppose \(S\) is bounded above, meaning
\(\exists\sup S = L\) as the least upper bound. Show that
\(\lim_{n \rightarrow \infty}a_{n} = L\). This means to show:

\begin{align*}
\forall\varepsilon &> 0,\ \exists N,\ n > N \Rightarrow \left| a_{n} - L \right| < \varepsilon \\
L - a_{n} &< \varepsilon\ \text{because\ we\ are\ proving\ the\ supremum}
\end{align*}


Let \(\varepsilon > 0\). We know:

\begin{align*}
&\exists a_{N} \in S\ \text{such\ that} \\
L - \varepsilon &< a_{N} \leq L
\end{align*}


Because it is between \(a_{n}\), which is greater than
\(L - \varepsilon\).

Show that \(\forall n > N,\ \ L - a_{n} < \varepsilon\)

We know \(N\) is fixed. If \(N = 100\), we need to show all \(n\) above
100. To prove this question, use induction.

\textbf{Base:} \(n = N + 1\). We're not trying to show for every natural
number \(n\) is true -- we're trying to show every natural number over
\(N\) is true.

\begin{itemize}
\item
  Show \(L - a_{N + 1} < \varepsilon\). Our assumption is we know
  \(L - \varepsilon < a_{N} \leq L\) (within epsilon distance). But
  \(a_{N + 1}\) should be the right (or the same position) as \(a_{N}\).
  Manipulate the inequalities a bit:
\item
  \(a_{N + 1} \geq a_{N}\)
\item
  \(- a_{N + 1} \leq - a_{N}\) (Multiply both sides by \(- 1\))
\item
  \(L - a_{N + 1} \leq L - a_{N}\)
\end{itemize}

Remember \(L - \varepsilon < a_{N} \leq L\) (which we assumed true)?
Rearrange this:

\begin{align*}
L - \varepsilon &< a_{N} \leq L \\
- \varepsilon &< a_{N} - L \leq 0 \\
\varepsilon &> L - a_{N}
\end{align*}


So:

\begin{itemize}
\item
  \(L - a_{N + 1} \leq L - a_{N} < \varepsilon\)
\end{itemize}

We showed what we wanted to show.

\textbf{Inductive step:}

\begin{itemize}
\item
  Suppose \(L - a_{n} < \varepsilon\)
\item
  Show \(L - a_{n + 1} < \varepsilon\)
\end{itemize}

We know

\begin{itemize}
\item
  \(a_{n + 1} \geq a_{n}\) because the sequence is increasing
\item
  \(- a_{n + 1} \leq - a_{n}\)
\item
  \(L - a_{n + 1} \leq L - a_{n}\)
\item
  \(L - a_{n + 1} \leq L - a_{n} < \varepsilon\) by the inductive
  hypothesis.
\end{itemize}

Since \(L = \sup S\), \(a_{n} \leq L\), so
\(\left| a_{n} - L \right| = L - a_{n}\)


\end{proof}

What is \(a_{N}\)? The thing you bump into when taking a small step to
the left from the supremum. Then you will bump into something from
\(S\). Then you can assume it is greater than \(L - \varepsilon\)
because the step is very small.

Flip this definition for decreasing equivalents. If \(S\) is bounded
below and it's \textbf{non-increasing,} then it converges to its
\textbf{infimum.}

Being \textbf{non-increasing} or \textbf{non-decreasing} is called
monotone.

\emph{The limit of a sequence has nothing to do with how it behaves at
the beginning.} Feel free to ignore the first few terms (or more,
precisely finitely as many).


\begin{figure}[H]
\centering
\includegraphics[width=4.0839in,height=0.55216in]{latex_images_Sequences_20/media/image1.png}
\end{figure}



\subsection{Finding limits to
infinity}

Given in a regular formula, compute the limit to infinity like normal.
For example:

\[\frac{n + ( - 1)^{n}}{n} = \frac{n}{n} + \frac{( - 1)^{n}}{n} = 1 + 0 = 1\]

And

\[{1.01}^{n} = \left( e^{\ln(1.01)} \right)^{n} = e^{n \cdot \ln(1.01)} = e^{\infty} = \infty\]

And

\[\frac{4n}{\sqrt{4n^{2} + 2}} = \frac{\frac{4n}{n}}{\sqrt{\frac{4n^{2}}{n^{2}} + \frac{2}{n^{2}}}} = \frac{4}{\sqrt{4}} = \frac{4}{2} = 2\]


\subsection{Recursively defined
sequences}

Consider the recursively defined sequence. Determine whether it
converges. If it converges, find its limit.

\begin{align*}
a_{1} &= 1, \\
a_{n + 1} &= \frac{1}{2}a_{n} + 1
\end{align*}


This means:

\begin{align*}
a_{2} &= \frac{1}{2} + 1 = \frac{3}{2} \\
a_{3} &= \frac{3}{4} + 1 = \frac{7}{4} \\
a_{4} &= \frac{7}{8} + 1 = \frac{15}{8}
\end{align*}


How can we show that it converges? We need to show that this function
is:

\begin{itemize}
\item
  Bounded above
\item
  Increasing / non-decreasing
\item
  And find \(\sup S = L\) (the supremum)
\end{itemize}

So, what do we do? Take the limit of both sides in the recursive part of
the function.

\begin{align*}
\lim_{n \rightarrow \infty}a_{n + 1} &= \lim_{n \rightarrow \infty}{\frac{1}{2}a_{n} + 1} \\
L &= \frac{1}{2}L + 1 \\
L &= 2
\end{align*}


Here, we kind of assumed that \(\infty + 1 = \infty\), but assuming that
isn't sufficient. We need to show that it is bounded above, so we're
going to use induction. We technically aren't supposed to do that first,
but we're going to do it anyways, because we need to so we can prove it.

Show bounded above: show \(\forall n,\ a_{n} \leq 2\).

\begin{proof} Base: \(n = 1,\ a_{1} = 1 \leq 2\)

Suppose \(a_{n} \leq 2\). Show \(a_{n + 1} \leq 2\)

\[a_{n + 1} = \frac{1}{2}a_{n} + 1 \leq \frac{1}{2}2 + 1 = 1 + 1 = 2\ \]
\end{proof}


\subsubsection{How to show its
increasing}

For \(a_{1} = 1,\ a_{n + 1} = \frac{1}{2}a_{n} + 1\)

\begin{itemize}
\item
  Show \(\forall n,\ a_{n + 1} > a_{n}\). Use induction:
\item
  Base: \(n = 1\). In this case, \(a_{2} = \frac{3}{2},\ a_{1} = 1\).
  Then \(a_{2} > a_{1}\)
\item
  Suppose \(a_{n + 1} > a_{n}\). Then show \(a_{n + 2} > a_{n + 1}\).
\item
  Because we know \(a_{n + 1} > a_{n}\),
\item
  \(a_{n + 2} = \frac{1}{2}a_{n + 1} + 1 > \frac{1}{2}a_{n} + 1 = a_{n + 1}\)
\end{itemize}


\subsection{The big theorem}

Constant \textless{} Logs \textless{} Polynomials \textless{}
Exponentials \textless{} Factorials \textless{} Power of itself

We get the ascending growth rates:

\[1 \ll \ln(n) \ll \sqrt{n} \ll n \ll n^{2} \ll \cdots \ll 2^{n} \ll e^{n} \ll n! \ll n^{b}\]

Where \(a_{n} \ll b_{n}\),

\[\lim_{n \rightarrow \infty}\frac{a_{n}}{b_{n}} = 0\]


\section{Improper integrals}

The integral is the area of many rectangles, which is base times height.

\textbf{Type 1: Base is infinite}

\[\int_{a}^{\infty}{f(x)\text{d}x} = \lim_{b \rightarrow \infty}{\int_{a}^{b}{\text{\ f}(x)\text{d}x}}\]

\textbf{Type 2: Height is infinite}

The function is approaching infinity. Assume \(f(x)\) approaches
infinity at \(x = k\) and \(a < k\). Then:

\[\int_{a}^{k}{\text{\ f}(x)} = \lim_{b \rightarrow k^{-}}{\int_{a}^{b}{\text{\ f}(x)\text{d}x}}\]

Example: The harmonic series:

\[\int_{0}^{1}{\frac{1}{x}\text{d}x} = \ln|x| \mid_{0}^{1} = \ln 1 - \ln 0 = 0 - ( - \infty) = \infty\]

Alternatively:

\[\lim_{b \rightarrow 0^{+}}{\ln|x|} \mid_{b}^{1} = 0 - ( - \infty) = \infty\]

Some more problems regarding improper integrals:

\[\int_{- 1}^{\infty}\frac{1}{x^{2} + 1} = \arctan(x) \mid_{- 1}^{\infty} \rightarrow \lim_{b \rightarrow \infty}{\arctan(x) \mid_{- 1}^{b}\ }\]

\[= \arctan(\infty) - \arctan( - 1) = \frac{\pi}{2} - \left( - \frac{\pi}{4} \right) = \frac{3\pi}{4}\]

Example 2:

\[\int_{2}^{\infty}{\frac{1}{x + 2}\text{d}x} = \ln|x + 2| \mid_{2}^{\infty} = \infty\]

Example 3:

\[\int_{0}^{\infty}{\cos(x)\text{d}x} = \sin(x) \mid_{0}^{\infty} \rightarrow \lim_{b \rightarrow \infty}{\sin(x) \mid_{0}^{b}} = \text{DNE}\]

Example 4:

\[\int_{0}^{1}{\frac{1}{x^{2}}\text{d}x} = \int_{0}^{1}x^{- 2} = \frac{x^{- 1}}{- 1} \mid_{0}^{1} = - \frac{1}{x} \mid_{0}^{1} = - \frac{1}{1} - \left( - \frac{1}{0} \right) = + \infty\]

Of course, we should have used
\(\lim_{b \rightarrow 0^{+}}{\frac{1}{x} \mid_{b}^{1}} = \infty\)

Example 5:

\[\int_{0}^{1}{\frac{1}{\sqrt{x}}\text{d}x} = \int_{}^{}x^{- \frac{1}{2}} = \frac{x^{\frac{1}{2}}}{\frac{1}{2}} = 2\sqrt{x} \mid_{0}^{1} = 2 - 0 = 2\ \]

It turns out that somehow when the function is \(\frac{1}{\sqrt{x}}\).
The area is finite. You can't see that the area is finite, but it is
finite. Of course, we should have used
\(\lim_{b \rightarrow 0^{+}}{2\sqrt{x} \mid_{b}^{1}}\), but that's just
a technical detail.


\subsection{Multiple infinite points}

Split the integral into several pieces, so that in each piece, only 1 of
the 2 bounds is approaching an infinite point. For example:

\[\int_{0}^{\infty}\frac{1}{(x - 1)^{3}}\]

Just integrate from 0 to 1, 1 to some finite location \(C\), and \(C\)
to infinity.


\begin{figure}[H]
\centering
\includegraphics[width=4.5in,height=1.10069in]{latex_images_Sequences_20/media/image2.png}
\end{figure}


This changes the direction of limits if you choose to evaluate some. In
this case:

\begin{align*}
&\lim_{b \rightarrow 1^{-}}{\int_{0}^{b}{\frac{1}{(x - 1)^{3}}\text{d}x}} + \lim_{b \rightarrow 1^{+}}{\int_{b}^{c}\frac{1}{(x - 1)^{3}}\text{d}x}  \\
&+ \lim_{b \rightarrow \infty}{\int_{c}^{b}{\frac{1}{(x - 3)^{3}}\text{d}x}} \\
&= - \frac{2}{{2(x - 1)}^{2}} \mid_{0}^{1} + \frac{- 2}{{2(x - 1)}^{2}} \mid_{1}^{c} + \left( - \frac{2}{{2(x - 1)}^{2}} \mid_{c}^{\infty} \right)\  \\
&\cdots \\
&= - \infty + \infty + \frac{1}{2(c - 1)^{3}} = \pm \infty
\end{align*}


\textbf{When one piece of the integral results in infinity, the entire
integral diverges.}

The problem is you cannot cancel infinities. You need to split the
integral up if there are multiple infinite points.


\section{P-series}

Let \(p > 0\). For what values of \(p\) would the integral converge or
diverge?

\[\int_{1}^{\infty}{\frac{1}{x^{p}}\text{d}x}\]

Case 1: \(p = 1\). Then
\(\int_{}^{}\frac{1}{x} = \ln|x| \mid_{1}^{\infty} = \infty\)

\begin{align*}
&= \int_{}^{}x^{- p} = \frac{x^{- p + 1}}{- p + 1} \\
&\lim_{b \rightarrow \infty}b^{- p + 1}
\end{align*}


\[\lim_{x \rightarrow \infty}x^{n} = \infty,\]

What if \(n < 0\)? \(n = - k\). Then:

\[\lim_{x \rightarrow \infty}x^{- k} = \frac{1}{x^{k}} = 0\]

So, we get (\(b\) is like the upper bound of the integral, and avoids us
from typing in \(\infty\)):

\[\lim_{b \rightarrow \infty}b^{- p + 1} = \left\{ \begin{matrix}
\infty & p + 1 > 0\ \text{diverge} \\
0 & - p + 1 < 0\ \text{converge} \\
\end{matrix} \right.\ \ \]

So, when \(- p + 1 < 0\) is true, then we get a convergent interval. So:

\[\int_{1}^{\infty}\frac{1}{x^{p}}\]

Will converge if \(p > 1\) and will diverge if \(p \leq 1\).


\section{Integral comparison test}

Integrals can't always be computed, but we can ask whether an integral
is convergent.

Comparison test for improper integrals:

If \(f,\ g\) is continuous, and \(0 \leq f(x) \leq g(x)\) for
\(x \geq a\), then:

\[0 \leq \int_{a}^{\infty}{f(x)\text{d}x} \leq \int_{a}^{\infty}{g(x)\text{d}x}\]

In particular:

\begin{itemize}
\item
  To check for convergence, look for another integral that is larger and
  base it on that.
\item
  To check for divergence, look for another integral that is smaller
\end{itemize}

Example question:

\[\int_{1}^{\infty}\frac{\sin x + 2\cos(x) + 10}{x^{2}} \leq \int_{1}^{\infty}\frac{1 + 2 + 10}{x^{2}} < \infty\]


\subsection{Limit comparison test for improper
integrals}

If \(f,\ g\) are continuous and both test positive for \(x \geq a\), IF

\[\lim_{x \rightarrow \infty}\frac{f(x)}{g(x)} \neq 0 \neq \infty\]

Then:

\[f(x) \approx g(x)\ (\text{The\ rate\ of\ growth\ is\ similar)}\]

\[\int_{a}^{\infty}{f(x)}\text{dx\ }\text{converges\ if\ and\ only\ if}\ \int_{a}^{\infty}{g(x)\text{d}x}\ \text{converges}\]

For example:

\[\int_{}^{}\frac{1}{x^{2}} < \infty,\ \int_{}^{}{\frac{1}{2x^{2}} < \infty},\ \int_{}^{}\frac{1}{x^{2} + 1} < \infty\]

Checking limits:

\[\lim\left( \frac{f}{g} \right) = \lim\left( \frac{1}{x^{2}}\left( x^{2} + 1 \right) \right) = 1\]

Hence why we don't want \(f\) to grow faster than \(g\) and \(g\) to
grow faster than \(f\). Similar to them having the same \(\Theta\).

Example:

\[\int_{0}^{\infty}{\frac{(x - 7)}{x^{2} + x + 5}\text{d}x} \approx \int_{1}^{\infty}\frac{x}{x^{2}} = \int_{}^{}\frac{1}{x} = \infty\]

Does \(x - 7\) and \(x^{2} + x + 5\) grow at the same rate?

\[\frac{(x - 7)}{x^{2} + x + 5}\]

Then we do some:

\[\frac{(x - 7)}{x^{2} + x + 5} \cdot \frac{x^{2}}{x}\]


\section{Series}

Consider the sequence
\(s_{1} = a_{1},\ s_{2} = a_{1} + a_{2},\ s_{3} = a_{1} + a_{2} + a_{3},\ \cdots\)

The sequence \(\{ s_{n}\}\) are the partial sums. We say that the series
converges if the sequence \(\{ s_{n}\}\) converges.

Usually, infinite series are very difficult to calculate, but for
geometric series:

\[\sum_{x = 0}^{\infty}x^{n} = \frac{1}{1 - x},\ |x| < 1\]

Only when series converge, we can perform regular algebra with them like
non-infinite series.

\begin{proof}

\[\sum_{x = 0}^{\infty}x^{n} = \lim_{N \rightarrow \infty}{\sum_{n = 0}^{N}x^{n}}\]

First step: \(\sum_{n = 0}^{N}x^{n} = \frac{1 - x^{N + 1}}{1 - x}\).

Then as \(N \rightarrow \infty\), \(x^{N + 1} = 0\), as \(|x| < 1\).
Then we get \(\frac{1}{1 - x}\).

\begin{align*}
&\sum_{n = 0}^{N}x^{n} = \frac{1 - x^{N + 1}}{1 - x} = 1 + x + x^{2} + \ldots + x^{N} \\
(1 - x)\left( 1 + x + x^{2} + \cdots + x^{N} \right) &= 1 - x^{N + 1}
\end{align*}


If we expand the left side:

\[1 + x + x^{2} + \cdots + x^{N} - x - x^{2} - \cdots - x^{N} - x^{N + 1} = 1 - x^{N + 1}\]

Now that we have the formula, we can then figure out the sum of infinite
series.


\end{proof}


\subsection{When series converge}

We can ignore finitely as many terms in the front of the sequence.


\subsubsection{The zero test for
divergence}

If \(\lim_{n \rightarrow \infty}a_{n} \neq 0\), then
\(\sum_{n = 1}^{\infty}a_{n}\) diverges. This test is almost never
useful due to oscillations or divergence for other reasons.


\subsubsection{The integral test}

For \(a_{n}\) positive and decreasing eventually:

\[\sum_{n = 1}^{\infty}a_{n} \approx \int_{1}^{\infty}a_{n}\]

The series converges if and only if the integral converges. Almost never
useful due to how hard it is to calculate integrals.

\textbf{Harmonic series}

\begin{align*}
&\sum_{n = 1}^{\infty}\frac{1}{n} = 1 + \frac{1}{2} + \frac{1}{3} + \frac{1}{4} + \cdots = \infty \\
\int_{1}^{\infty}{\frac{1}{x}\text{d}x} &= \ln|x| \mid_{1}^{\infty} = \ln\infty
\end{align*}



\subsubsection{P-series}

\[\sum_{}^{\infty}\frac{1}{n^{p}} \approx \int_{}^{\infty}{\frac{1}{x^{p}}\text{d}x}\]

Converges for \(p > 1\), diverges for \(p \leq 1\)

Given this result, what can we say about:

\[\int_{}^{\infty}\frac{1}{n^{2} - 1} \approx \sum_{}^{\infty}\frac{1}{n^{2}}\]

This will converge as \(p > 1\). When \(a_{n}\) is a rational function,
we only need to take the leading powers on the top and bottom to create
a new series \(\sum_{}^{\infty}b_{n}\).


\subsubsection{Limit comparison test}

For \(a_{n},\ b_{n}\), positive eventually:

\[\lim_{n \rightarrow \infty}\frac{a_{n}}{b_{n}} \neq 0 \neq \infty \Rightarrow \sum_{}^{\infty}a_{n} \approx \sum_{}^{\infty}b_{n}\]

Example: \(a_{n} = n^{2},\ b_{n} = n^{2} - 1\)

\[\frac{a_{n}}{b_{n}} = \frac{1}{n^{2} - 1} \cdot \frac{n^{2}}{1} = 1\]


\subsubsection{Comparison tests}

Suppose \(a_{n} \leq b_{n}\) and positive eventually.

If \(\sum_{}^{\infty}b_{n}\) converges then \(\sum_{}^{\infty}a_{n}\)
converges.

If \(\sum_{}^{\infty}b_{n}\) diverges then \(\sum_{}^{\infty}a_{n}\)
diverges\ldots?

For example:

\[\sum_{}^{\infty}\frac{\arctan(n)}{n^{2}} \leq \sum_{}^{\infty}{\frac{\pi}{2} \cdot \frac{1}{n^{2}}}\]

This converges.

Using \(\ln(n) < n^{p}\) for any \(p > 0\) (pick \(p\) to be small):

\[\sum_{}^{\infty}\frac{\ln(n)}{n^{2}} < \sum_{}^{\infty}\frac{\sqrt{n}}{n^{2}} = \sum_{}^{\infty}\frac{1}{n^{\frac{3}{2}}}\]

Which converges.


\subsubsection{Ratio test}

For \(a_{n}\) positive eventually, compute:

\[\lim_{n \rightarrow \infty}\frac{a_{n + 1}}{a_{n}} = L\]

If \(L < 1\), then \(\sum_{}^{\infty}a_{n}\) converges.

If \(L > 1\), then \(\sum_{}^{\infty}a_{n}\) diverges.

If \(L = 1\), it is inclusive.

Example:

\begin{align*}
\frac{a_{n + 1}}{a_{n}} &= L < 1 \Rightarrow  \\
&\lim_{n \rightarrow \infty}{\left( \frac{1}{2} \right)^{n + 1} \cdot \frac{2^{n}}{1} = \frac{1}{2}} \\
a_{n + 1} &\approx \frac{1}{2}a_{n}
\end{align*}


\(L = 1\) example:

\[\frac{a_{n + 1}}{a_{n}} = \frac{1}{n + 1} \cdot \frac{n}{1} = 1\]

Not a huge difference between \(a_{n + 1}\) and \(a_{n}\). So, this test
does not give us any conclusions, as we don't know how fast this is
decreasing. This means this test does not work for polynomials, but it
does for exponentials and factorials. For polynomials, use the p-series.

\(L > 1\) example: If \(\frac{a_{n + 1}}{a_{n}} > 1\), it's getting
larger and larger. The series is divergent.


\subsection{Exponential vs. Polynomial
series}

Exponential series can be calculated easily. Polynomials, on the other
hand, seem to be very hard.

\[\sum_{n = 1}^{\infty}\frac{1}{n^{2}} = \frac{\pi^{2}}{6}\]


\subsection{Convergence of series involving positive or negative
terms}

\[\sum_{}^{}a_{n},\ a_{n}\sim\frac{1}{n} = \frac{1}{1} - \frac{1}{2} - \frac{1}{3} + \frac{1}{4} - \frac{1}{5} + \frac{1}{6} + \frac{1}{7} + \frac{1}{8} - \frac{1}{9} + \frac{1}{10} - \frac{1}{11} \pm \cdots\]

In general, this is a difficult problem. Thus, it is much easier to just
get rid of all the negative signs. Now, with this series, the negative
terms are somehow cancelling some of the positive terms. This means it
is possible for this series to not be infinite. We've said that:

\[\sum_{}^{}a_{n} \leq \sum_{}^{}\left| a_{n} \right|\]

There may be cancellation on the left side, for our above series. But
there is no cancellation on the right.

\textbf{Test 5: Absolute Convergence Test}

\[if\sum_{}^{\infty}\left| a_{n} \right|\ \text{converges} \Rightarrow \text{then}\sum_{}^{\infty}a_{n}\ \text{converges}\]

When \(\sum_{}^{\infty}\left| a_{n} \right|\) converges, we say
\(\sum_{}^{\infty}a_{n}\) is \textbf{absolutely convergent.} Also, if
\(a_{n} > 0\ \forall n,\ \)convergent and absolutely convergent mean the
same thing.

There is a distinction when \(\exists a_{n},\ a_{n} < 0\). In the event
where the original series involving positive and negative terms
converge, but if we absolute value \(a_{n}\), then we say that it is
conditionally convergent. In formal logic:

\[\sum_{}^{\infty}a_{n}\ \text{converges}\ \text{and\ }\sum_{}^{\infty}\left| a_{n} \right|\ \text{diverges} \Rightarrow \sum_{}^{\infty}a_{n}\ \text{is\ conditionally\ convergent}\]

An example of an erratic sequence:

\[\sum_{}^{\infty}\left| \frac{\sin(n)}{n^{2}} \right| \leq \sum_{}^{\infty}\frac{1}{n^{2}}\ \text{which\ converges}\]

The sine function is erratic, as its period is \(2\pi\), which happens
to not be a rational number or even an integer. Determining a pattern
for \(\sin(n)\) being positive or negative with increasing \(n\)s can be
very hard to determine by hand, so the distributions of + and \(-\)s are
very erratic.

Well:

\begin{align*}
\sin(n) &\leq 1 \Rightarrow  \\
\sum_{}^{\infty}\left| \frac{\sin(n)}{n^{2}} \right| &\leq \sum_{}^{\infty}\frac{1}{n^{2}} \Rightarrow \sum_{}^{\infty}\frac{\sin(n)}{n^{2}}\ \text{converges\ absolutely}
\end{align*}


By the theorem we stated in the beginning of this subsection.


\subsubsection{The alternating series
test}

A special case is when the positive and negative terms alternate in
turn:

\[a_{1} - a_{2} + a_{3} - a_{4}\ldots\]

This can be written with the notation \(( - 1)^{n}\) or
\(( - 1)^{n + 1}\).

For example:

\[\sum( - 1)^{n} \times \frac{1}{n}\sim\frac{1}{1} - \frac{1}{2} + \frac{1}{3} - \frac{1}{4} + \ldots\]

We know:

\[\sum_{}^{\infty}\left| \frac{( - 1)^{n}}{n} \right| = \sum_{}^{\infty}\frac{1}{n} = \infty\]

(When you have infinitely as many things, you cannot change the order in
which you add terms).

And also, we have this theorem, which requires us to only compute the
limit of \(a_{n}\):

\[\lim_{(n \rightarrow \infty)}a_{n} = 0\ \text{if\ and\ only\ if}\ \sum_{}^{\infty}{( - 1)^{n}a_{n}}\ \text{converges}\]

In the event, when you have positive and negative terms alternating, we
don't need to worry how fast the terms head towards zero. Essentially
what happens, is that even if 0 is being approached slowly, with
adjacent positive and negative terms, we still get convergence.

\begin{align*}
&\left( \frac{1}{1} - \frac{1}{2} \right) + \left( \frac{1}{3} - \frac{1}{4} \right) + \ldots \\
&= \frac{1}{2} + \frac{1}{12} + \ldots
\end{align*}


Because \(\lim_{n \rightarrow \infty}\frac{1}{n} = 0\) we know that
\(\sum_{}^{\infty}\frac{( - 1)^{n}}{n}\) converges, using the theorem
from above.

The test also holds with
\(\cos\left( n\pi  \right) = ( - 1)^{n} = \sin\left( \left( n + \frac{1}{2} \right)\pi \right) = ( - 1)^{n + 1}\)


\paragraph{A note about looping
convergence}

\(a_{1} - a_{2} + a_{3} - a_{4}\).

Add \(a_{1}\) to 0. Then subtract \(a_{2}\). But since you said the
sequence is decreasing:

\[a_{1} > a_{1} - a_{2} > 0\]

And then you add back \(a_{3}\). Then you take back most of what you
subtracted back in \(a_{2}\). So:

\[a_{1} > a_{1} - a_{2} + a_{3} > a_{1} - a_{2} > 0\]

The cycle repeats.


\begin{figure}[H]
\centering
\includegraphics[width=2.18518in,height=2.22363in]{latex_images_Sequences_20/media/image3.png}
\end{figure}


And that's another reason to why \(\sum_{}^{\infty}\frac{1}{n}\)
converges.


\section{Summary of Tests for
Series}

Terms lower under \(+\) can be ignored, but terms under multiplication
can't be ignored. You can say that the numerator / denominator is order
something\emph{.}

Combining all the tests, the steps for determining convergence of a
series \(\sum_{}^{\infty}a_{n}\):

\begin{enumerate}
\def\labelenumi{\arabic{enumi}.}
\item
  Take absolute value of \(a_{n}\): \(\left| a_{n} \right|\)
\item
  If there are exponents or factorials, apply the ratio test:
  \(\lim_{n \rightarrow \infty}\left| \frac{a_{n + 1}}{a_{n}} \right| = L\)
  or the root test:
  \(\lim_{n \rightarrow \infty}\sqrt[n]{\left| a_{n} \right|}\)

  \begin{enumerate}
  \def\labelenumii{\alph{enumii}.}
  \item
    \(L < 1\) concludes \(\sum_{}^{\infty}a_{n}\) absolutely converges.
  \item
    \(L > 1\) concludes \(\sum_{}^{\infty}a_{n}\) diverges.
  \item
    \(L = 1\) and we are unclear.
  \end{enumerate}
\item
  Upper bounds:

  \begin{enumerate}
  \def\labelenumii{\alph{enumii}.}
  \item
    If \(\ln(n)\) in the numerator, compare \(1 < \ln(n) < n^{p}\) with
    \(p\) small.
  \item
    \(\ln(n)\) in the denominator, compare
    \(1 > \frac{1}{\ln(n)} > \frac{1}{n^{p}}\) with \(p\) small.
  \item
    For sine and cosine in the numerator, conclude
    \(\left| \sin(n) \right| \leq 1,\ |cos(n)| \leq 1\).
  \item
    For arctangent in the numerator, compare
    \(\arctan(n) < \frac{\pi}{2}\)
  \end{enumerate}
\item
  For series in rational function form, take \(b_{n}\) to be leading
  powers of top and bottom. Attain \(b_{n} = \frac{1}{n^{p}}\).

  \begin{enumerate}
  \def\labelenumii{\alph{enumii}.}
  \item
    \(p > 1 \Rightarrow \sum_{}^{\infty}a_{n}\) absolutely converges.
  \item
    \(p \leq 1 \cdot \sum_{}^{\infty}{|a_{n}|}\) diverges, but
    \(\sum_{}^{\infty}a_{n}\) may still converge.
  \end{enumerate}
\item
  If the original series can be written in the form
  \(\sum_{}^{\infty}{( - 1)^{n}a_{n}}\), compute
  \(\lim_{n \rightarrow \infty}a_{n} = L\).

  \begin{enumerate}
  \def\labelenumii{\alph{enumii}.}
  \item
    \(L = 0 \Rightarrow \sum_{}^{\infty}{( - 1)^{n}a_{n}}\)
    conditionally converges.
  \item
    \(L \neq 0 \Rightarrow \sum_{}^{\infty}{( - 1)^{n}a_{n}}\) diverges.
  \end{enumerate}
\end{enumerate}

\[\sum_{}^{\infty}a_{n}\ \text{converges}\ \text{and\ }\sum_{}^{\infty}\left| a_{n} \right|\ \text{diverges} \Rightarrow \sum_{}^{\infty}a_{n}\ \text{is\ conditionally\ convergent}\]

Otherwise, it is absolutely convergent.


\subsection{Linearithmic on the
denominator}

For \(a_{n} = \frac{1}{n \cdot \ln(n)}\), there is one power of \(n\) in
the bottom so use the integral test:
\(\int_{}^{}\frac{1}{x\ln(x)}\text{d}x\)

Use small angle approximation on (note that
\(n^{p} \rightarrow \infty \Rightarrow \frac{1}{n^{p}} \rightarrow 0\)):

\[\sin\left( \frac{1}{n^{p}} \right) \approx \frac{1}{n^{p}} \approx \tan\left( \frac{1}{n^{p}} \right) \approx \arcsin\left( \frac{1}{n^{p}} \right) \approx \arctan\left( \frac{1}{n^{p}} \right)\]


\subsection{The very obvious}

If
\(\lim_{n \rightarrow \infty}a_{n} \neq 0 \Rightarrow \sum_{}^{\infty}a_{n}\)
diverges.


\subsection{Some examples}

Special case 1: From something you know, try to get to the original
function.

\begin{align*}
&\sum_{n = 1}^{\infty}\frac{1}{2 + \cos(n)} \\
- 1 &\leq \cos(n) \leq 1 \\
1 &\leq 2 + \cos(n) \leq 3 \\
\frac{1}{3} &\leq \frac{1}{2 + \cos(n)} \leq 1 \\
\lim_{n \rightarrow \infty}\frac{1}{2 + \cos(n)} &\neq 0
\end{align*}


Divergent.

Special case 2:

\begin{align*}
\sum_{n = 2}^{\infty}\frac{1}{n\ln(n)} &\approx \int_{2}^{\infty}{\frac{1}{x\ln(x)}\text{d}x} \\
u &= \ln(x),\ du = \frac{1}{x}\text{d}x \\
\int_{}^{}{\frac{1}{u}\text{d}u} &= \ln|u| = \ln\left| \ln(x) \right| \mid_{a}^{\infty} \\
&= \lim_{b \rightarrow \infty}{\ln\left| \ln b \right|} = \infty
\end{align*}


Divergent.

(3)

\begin{align*}
&\sum_{n = 1}^{\infty}\frac{( - 1)^{n}}{n} \\
&\sum_{n = 1}^{\infty}\left| \frac{( - 1)^{n}}{n} \right| = \sum_{}^{}\frac{1}{n}
\end{align*}


Because \(p = 1 \leq 1\), this one diverges.

Consider \(\lim_{n \rightarrow \infty}\frac{1}{n} = 0\). Then we can
tell the series of (3) is conditionally convergent. Conditionally
converges \(\Rightarrow\) converges.

(4)

\[\sum_{n = 1}^{\infty}\frac{1}{n^{2}}\]

Converges due to the \(p\)-series

(5)

\[\sum_{n = 1}^{\infty}\frac{( - 1)^{n}}{n^{2} - 1},\ \sum_{n = 1}^{\infty}\left| \frac{( - 1)^{n}}{n^{2} - 1} \right| \approx \sum_{n = 1}^{\infty}\left| \frac{1}{n^{2}} \right| = \sum_{}^{}\frac{1}{n^{2}}\]

This absolutely converges. Meaning it converges.

(6)

\begin{align*}
&\sum_{n = 1}^{\infty}\frac{n^{2} + 4n}{\sqrt{n^{5} + 4n^{4} + 2n}} \\
\sum_{n = 1}^{\infty}\left| \frac{n^{2} + 4n}{\sqrt{n^{5} + 4n^{4} + 2n}} \right| &\approx \sum_{n = 1}^{\infty}\left| \frac{n^{2}}{\sqrt{n^{5}}} \right| = \sum_{n = 1}^{\infty}\frac{1}{n^{\frac{1}{2}}},\ p = \frac{1}{2} \leq 1
\end{align*}


Diverges due to the p-series.

(7)

\begin{align*}
&\sum_{n = 1}^{\infty}\frac{\sqrt{n}\sin(n)}{n^{2}} \\
\sum_{n = 1}^{\infty}\left| \frac{\sqrt{n}\sin(n)}{n^{2}} \right| &\leq \sum_{n = 1}^{\infty}\left| \frac{\sqrt{n}}{n^{2}} \right| = \sum_{n = 1}^{\infty}\left| \frac{1}{n^{\frac{3}{2}}} \right|
\end{align*}


Converges due to the p-series. This converges absolutely.

(8)

\begin{align*}
&\sum_{n = 1}^{\infty}{( - 1)^{n}\sin\left( \frac{1}{n} \right)} \\
\sum_{n = 1}^{\infty}\left| ( - 1)^{n}\sin\left( \frac{1}{n} \right) \right| = \sum_{n = 1}^{\infty}{\sin\left( \frac{1}{n} \right)} &\approx \sum_{n = 1}^{\infty}\frac{1}{n}
\end{align*}


Divergent due to the p-series. But does it conditionally converge?

\[\lim_{n \rightarrow \infty}{\sin\left( \frac{1}{n} \right)} = 0\ \text{converges}\]

So, it conditionally converges.

(9)

\[\sum_{}^{}\left| ( - 1)^{n}\sin\left( \frac{1}{n^{2}} \right) \right| = \sum_{}^{}{\sin\left( \frac{1}{n^{2}} \right)} \approx \sum_{}^{}\frac{1}{n^{2}}\]

Converges absolutely.


\subsubsection{An aside on small angle
approximation}

\begin{align*}
\sum_{n = 1}^{\infty}{\sin\left( \frac{1}{n^{2}} \right)} &\approx \sum_{n = 1}^{\infty}\frac{1}{n^{2}} \\
\lim_{n \rightarrow \infty}{\sin\left( \frac{1}{n^{2}} \right) \cdot n^{2}} &\approx \frac{1}{n^{2}} \cdot n^{2}
\end{align*}


(10)

\[\sum_{n = 1}^{\infty}\frac{(n + 1)2^{n}}{n!}\]

Time for the ratio test:

\begin{align*}
\lim_{n \rightarrow \infty\ }\left| \frac{a_{n + 1}}{a_{n}} \right| &= \left| \frac{(n + 2)2^{n + 1}}{(n + 1)!} \cdot \frac{n!}{(n + 1)2^{n}} \right| = \left( \frac{n!}{(n + 1)!} \cdot \frac{n + 2}{n + 1} \cdot \frac{2^{n + 1}}{2^{n}} \right) \\
&= \left( \frac{1}{n + 1} \cdot \frac{n + 2}{n + 1} \cdot \frac{2^{n + 1}}{2^{n}} \right) = 0 \cdot 1 \cdot 2 = 0
\end{align*}


Meaning this converges absolutely.

The point of the ratio test is to take the next term and divide it by
the previous term. Then, you can say that the series decreases fast
enough.

(11)

\[\sum_{n = 1}^{\infty}\frac{n!}{2^{3n}}\]

Try

\begin{align*}
\lim_{n \rightarrow \infty}{\ \left| \frac{a_{n + 1}}{a_{n}} \right|} &= \left| \frac{(n + 1)!}{2^{3n + 3}} \cdot \frac{2^{3n}}{n!} \right| \\
&= \left| \frac{(n + 1)}{2^{3}} \right|
\end{align*}


Diverges

(12)

\[\sum_{n = 1}^{\infty}\frac{n!}{(2n)!}\]

\begin{align*}
&\lim_{n \rightarrow \infty}\left| \frac{a_{n + 1}}{a_{n}} \right| \\
&= \frac{(n + 1)!}{\left( 2(n + 1) \right)!} \cdot \frac{(2n)!}{n!} = \frac{(n + 1)n!}{(2n + 2)(2n + 1)(2n)!} \cdot \frac{(2n)!}{n!} = \frac{n + 1}{(2n + 2)(2n + 1)} \approx \frac{n}{n^{2}} \approx \frac{1}{n}
\end{align*}


Diverges

(13)

\[\sum_{n = 1}^{\infty}\frac{n!(2n)!}{(3n)!}\]

The ratio test:

\begin{align*}
&\frac{(n + 1)!(2n + 2)!}{(3n + 3)!} \cdot \frac{(3n)!}{n!(2n)!} \\
&= \frac{(3n)!}{(3n + 3)!} \cdot \frac{(n + 1)!}{n!} \cdot \frac{(2n + 2)!}{(2n)!} \\
&= \frac{(3n)!}{(3n + 3)(3n + 2)(3n + 1)(3n)!} \cdot \frac{(n + 1)n!}{n!} \cdot \frac{(2n + 2)(2n + 1)(2n)!}{(2n)!} \\
&= \frac{1}{(3n + 3)(3n + 2)(3n + 1)} \cdot \frac{(n + 1)}{1} \cdot \frac{(2n + 2)(2n + 1)}{1}
\end{align*}


There are 3 powers of the numerator and 3 powers on the denominator.

Expand:

\[\frac{\left( 2n^{2} + 4n + 2 \right)(2n + 1)}{\left( 9n^{2} + 15n + 6 \right)(3n + 1)} = \frac{4n^{3} + \cdots}{27n^{3} + \cdots} = \frac{4}{27} < 1\]

It's convergent, absolutely. You can calculate this quickly by only
treating the leading coefficients as existing:

\[\frac{(1 \cdot 2 \cdot 2)n^{3}}{(3 \cdot 3 \cdot 3)n^{3}}\]

(14)

\[\sum_{n = 1}^{\infty}\frac{1}{n^{n}}\]

Let's try this:

\[\lim_{n \rightarrow \infty}{\frac{1}{(n + 1)^{n + 1}} \cdot n^{n}}\]

This becomes problematic. If you have \(n^{n}\) as your series, if you
try to use the ratio test, the answer may not be obvious. This means we
have to factor out \(n + 1\) in the exponent:

\[\frac{1}{(n + 1)^{n + 1}} \cdot \frac{1}{n + 1} = \left( \frac{n}{n + 1} \right)^{n} \cdot \frac{1}{n + 1} = \frac{1}{e} \cdot \frac{1}{n + 1} = 0 < 1\]

Okay, that isn't the best decision. However, since \(n^{n}\) sits on the
top of the growth hierarchy, we can do comparisons:

\[\sum_{n = 1}^{\infty}\frac{1}{n^{2}} > \sum_{n = 1}^{\infty}\frac{1}{n^{n}}\]

Because \(\sum_{n = 1}^{\infty}\frac{1}{n^{2}}\) converges, so does
\(\sum_{n = 1}^{\infty}\frac{1}{n^{n}}\). However, if we want to do the
ratio test:

Choose strategically, but make sure to choose something high in the
hierarchy that retains the information that \(n^{n}\) grows very fast.
For example:

\[\sum_{n = 1}^{\infty}\frac{1}{n^{n}}:\sum_{n = 1}^{\infty}\left| a_{n} \right| = \ \sum_{n = 1}^{\infty}\frac{2^{n}}{n^{n}}\sim\sum_{n = 1}^{\infty}\frac{2^{n}}{3^{n}},\ x = \frac{2}{3},\ |x| < 1 \Rightarrow \text{conver}\text{ges}\]

Try this again:

\begin{align*}
&\sum_{n = 1}^{\infty}\frac{1}{n^{n}} \\
&\lim_{n \rightarrow \infty}\left( \frac{n^{n}}{(n + 1)^{n + 1}} \right)
\end{align*}



\subsubsection{Evaluating n to the n}

From \(\frac{1}{n^{n}}\): Do the ratio test. We get
\(\left( \frac{n}{n + 1} \right)^{n} \cdot \frac{1}{n + 1}\). Also:
\(\frac{1}{(n + 1)^{n + 1}} \cdot n^{n} = \left( \frac{n}{n + 1} \right)^{n} \cdot \frac{1}{n + 1}\)

\begin{align*}
\lim_{n \rightarrow \infty}\left( \frac{n}{n + 1} \right)^{n} &= \frac{1}{e} \\
\left( \frac{n}{n + 1} \right) &= e^{\ln\left( \frac{n}{n + 1} \right)} \\
\left( \frac{n}{n + 1} \right)^{n} &= \left( e^{\ln\left( \frac{n}{n + 1} \right)} \right)^{n} = e^{n \cdot \ln\left( \frac{n}{n + 1} \right)}
\end{align*}


So what's
\(\lim_{n \rightarrow \infty}{n \cdot \ln\left( \frac{n}{n + 1} \right)}\)?
Perhaps use L'Hopital's rule. We have the form of something like
\(\infty \cdot 0\).

\[\frac{\ln\left( \frac{n}{n + 1} \right)}{\frac{1}{n}}\]

Take the derivative of both sides:

\[\frac{\frac{1}{\frac{n}{n + 1}} \cdot \frac{(n + 1) - n}{(n + 1)^{2}}}{- \frac{1}{n^{2}}} = - \frac{n + 1}{n} \cdot \frac{1}{(n + 1)^{2}} \cdot n^{2} = - \frac{1}{n} \cdot \frac{1}{n + 1} \cdot n^{2} = - \frac{n}{n + 1} = - 1\]

The answer to this exponent limit is \(- 1\) so the answer is
\(e^{- 1} = \frac{1}{e}\)


\subsubsection{The compound interest
limit}

\begin{align*}
\lim_{n \rightarrow \infty}\left( 1 + \frac{1}{n} \right)^{n} &= e \\
\left( \frac{n + 1}{n} \right)^{n} &= \left( \frac{n}{(n + 1)^{- 1}} \right)^{n} = \left( \left( \frac{n}{n + 1} \right)^{n} \right)^{- 1} = e
\end{align*}


Now, why is this \(e\)?
\(\lim_{n \rightarrow \infty}\left( 1 + \frac{1}{n} \right)^{n} = e\)

It is not defined directly. We have to define ln:

\[\ln(x) = \int_{1}^{x}{\frac{1}{t}\text{d}t}\]

But that was not how \(e\) was originally discovered. It was done
through this formula.

\[\lim_{n \rightarrow \infty}\left( 1 + \frac{1}{n} \right)^{n}\]

100\% interest? With \$1, one year later, you have \$1+\$1=\$2.

Now, you want a different plan. What if you withdraw mid-year? Well,
then I withdraw \$1.50. Keep it in the bank. \(\$ 1.50\). Half a year
later -- I was supposed to receive 100\% interest per year. So you get
\$1.5 + \$0.75. The total you now have is \$2.25. While previously you
had \$2 when the arrangement you had with the bank has always been the
same.

Say compounding happens \(\frac{1}{4}\) years. Then precisely at the
\(\frac{1}{4}\)\textsuperscript{th} point of the year, my money gets
multiplied by \(1\):

\begin{align*}
&\$ 1 \\
&\$ 1 + \$ 0.25 \\
\$ 1.25 + \frac{1}{4}\$ 1.25 &= \$ 1.25 + \$ 0.3125 =
\end{align*}


\[\lim_{n \rightarrow \infty}\left( 1 + \frac{1}{n} \right)^{n}\]

Must exist. So call it to be \(e\).

Why is this significant? These numbers are computationally infeasible at
high values.

Say we have \(r\%\) interest. For example, \(r = 0.05\).

Then \(\$ 1 + \$ 0.05 = 1 + \frac{r}{1}\)

So, for continuous compounding, we would have:

\[\lim_{n \rightarrow \infty}{\ \left( 1 + \frac{r}{n} \right)^{n}}\]

If you want an initial deposit:

\[P_{1} \approx \lim_{n \rightarrow \infty}{P_{0}\left( 1 + \frac{r}{n} \right)^{n}} = P_{0}e^{r}\]


\section{Taylor series}

To approximate the function \(f(x)\) with a polynomial.

We choose the point \(a\) as the center of expansion. We require the
polynomial to have the same derivatives as the function at \(a\). This
infinite polynomial is called the \textbf{Taylor (power) series of}
\(f(x)\)\textbf{.}

\[f(x) \approx P(x) = \sum_{n = 0}^{\infty}{\frac{f^{(n)}(a)}{n!}(x - a)^{n}}\]

Note that \(f'(a),\ f''(a),\ f^{(n)}(a)\) are constants. If
\(a = 0\), take the derivative and plug in 0.

\begin{align*}
P(x) &= f(a) + \frac{f'(a)}{1!}(x - a) + \frac{f''(a)}{2!}(x - a)^{2} + \frac{f'''(a)}{3!}(x - a)^{3} + \cdots \\
P'(x) &= f'(a) + f''(a)(x - a) + \frac{f'''(a)}{2!}(x - a)^{2} + \cdots \\
P''(x) &= f''(a)(1) + \frac{f'''(a)(x - a)}{2!} + \cdots \\
P'(a) &= f'(a) \\
P''(a) &= f''(a)
\end{align*}


We defined the derivative of the Taylor polynomial at a point to be
exactly the derivative of the original function\ldots{} at \(a\). The
same applies for the second derivative, and so on. The \(n!\) is here to
cancel out the powers that come out from \((x - a)^{n}\).

For \(a = 0\), the most important Taylor Series to know are:

\begin{align*}
e^{x} &= \sum_{n = 0}^{\infty}\frac{x^{n}}{n!} \\
\sin(x) &= \sum_{n = 0}^{\infty}{\frac{( - 1)^{n}}{(2n + 1)!}x^{2n + 1}} \\
\cos(x) &= \sum_{n = 0}^{\infty}{\frac{( - 1)^{n}}{(2n)!}x^{2n}} \\
\frac{1}{1 - x} &= \sum_{n = 0}^{\infty}x^{n} = 1 + x + x^{2} + x^{3} + \cdots,\ |x| < 1
\end{align*}


Usually, the approximation only holds in a symmetric interval centered
at \(a\). This is called the interval of convergence. Half the length of
the interval is called the radius of convergence.

It turns out, that the approximation for \(e^{x},\sin(x)\) and
\(\cos(x)\) works for every \(x\). Oftentimes, these functions are just
defined like these.

Since the formula holds for any \(x\) inside the interval of
convergence, we can substitute \(x\) with any other quantity.

\begin{align*}
e^{x^{2}} &= \sum_{n = 0}^{\infty}\frac{\left( x^{2} \right)^{n}}{n!} = 1 + x^{2} + \frac{x^{4}}{2!} + \frac{x^{6}}{3!} \\
\frac{1}{3 - x} &= \frac{1}{3} \cdot \frac{1}{1 - \frac{x}{3}} = \frac{1}{3}\sum_{n = 0}^{\infty}\left( \frac{x}{3} \right)^{n} = \frac{1}{3}\left( 1 + \frac{x}{3} + \left( \frac{x}{3} \right)^{2} + \left( \frac{x}{3} \right)^{3} + \cdots \right)
\end{align*}



\subsection{Derivatives and Integrals}

\begin{align*}
\int_{}^{}{\left( \sum a_{n} \right)\text{d}x} &= \sum_{}^{}{\int_{}^{}{a_{n}\text{d}x}} \\
\frac{d}{\text{d}x}\sum a_{n} &= \sum_{}^{}{\frac{d}{\text{d}x}a_{n}}
\end{align*}


Example:

\[\frac{1}{1 + x^{2}} = \frac{1}{1 - \left( - x^{2} \right)} = \sum_{n = 0}^{\infty}\left( - x^{2} \right)^{n} = \sum_{n = 0}^{\infty}{( - 1)^{n}x^{2n}}\]

Take the integral of both sides:

\begin{align*}
\int_{}^{}{\frac{1}{1 + x^{2}}\text{d}x} &= \arctan(x) = C + \sum_{n = 0}^{\infty}{\int_{}^{}{( - 1)^{n}x^{2n}}},\ |x| < 1 \\
\arctan(x) &= C + \sum_{n = 0}^{\infty}\frac{( - 1)^{n}x^{2n + 1}}{2n + 1} = \sum_{n = 0}^{\infty}\frac{( - 1)^{n}x^{2n + 1}}{2n + 1}
\end{align*}


This is the case where \(C\) is needed. To solve for \(C\), plug in
\(x = 0,\ x^{2n + 1} = 0\), so \(\arctan(0) = 0 = C + \sum_{}^{}0 = 0\)

Similarly:

\begin{align*}
\frac{1}{1 - x} &= \sum_{n = 0}^{\infty}x^{n} \\
\int_{}^{}{\frac{1}{1 - x}\text{d}x} &= C + \sum_{n = 0}^{\infty}{\int x^{n}\text{d}x},\ |x| < 1 \\
- \ln(1 - x) &= C + \sum_{n = 0}^{\infty}\frac{x^{n + 1}}{n + 1},\ |x| < 1 \\
- \ln(1 - x) &= C + \sum_{k = 1}^{\infty}\frac{x^{k}}{k}\ \text{to\ get\ rid\ of}\ n + 1,\ k = n + 1s \\
x &= 0,\  - \ln(1) = C + \sum_{n = 0}^{\infty}\frac{0^{n + 1}}{0 + 1} = 0,\ C = 0
\end{align*}


\[\ln(1 - x) = - \sum_{n = 0}^{\infty}\frac{x^{n}}{n} = - x - \frac{x^{2}}{2} - \frac{x^{3}}{3}\]


\subsection{Radius of convergence}

Involves \(n\) and \(x\) as a variable. Has an interval of convergence
and radius of convergence. For example, find the radius of convergence
for this in the power series:

\[\sum_{n = 1}^{\infty}{n(x - 2)^{n}}\]

Ratio test:

\begin{align*}
\lim_{n \rightarrow \infty}\left| \frac{a_{n + 1}}{a_{n}} \right| &= \lim_{n \rightarrow \infty}\left| \frac{(n + 1)(x - 2)^{n + 1}}{n(x - 2)^{n}} \right| = 1 \times |x - 2| \\
|x - 2| &< 1 \Rightarrow 1 < x < 3
\end{align*}


Another example:

\begin{align*}
&\sum_{n = 2}^{\infty}\frac{x^{3n}}{2^{n}} \\
\left| \frac{x^{3n + 3}}{2^{n + 1}} \cdot \frac{2^{n}}{x^{3n}} \right| &= \left| \frac{1}{2}x^{3} \right| < 1 \Rightarrow \left| x^{3} \right| < 2 \Rightarrow |x| < \sqrt[3]{2}
\end{align*}


Then 2 is the radius of convergence, and the interval of convergence is
\(- \sqrt[3]{2} < x < \sqrt[3]{2}\)

Isolate \(x\) when it is multiplied by a constant? Do it.

\begin{align*}
&\sum_{n = 3}^{\infty}{( - 1)^{n}(x - 3)^{2n}} \\
\lim_{n \rightarrow \infty}\left| \frac{a_{n + 1}}{a_{n}} \right| &= \left| \frac{( - 1)^{(n + 1)}(x - 3)^{2n + 2}}{( - 1)^{n}(x - 3)^{2n}} \right| \\
\left| (x - 3)^{2} \right| &< 1 \\
|x - 3| < \sqrt{1} &= 1 \\
2 &< x < 4
\end{align*}



\subsection{Taylor series of some
functions:}

\[\cos\left( x^{4} \right) = 1 - \frac{x^{8}}{2!} + \frac{x^{16}}{4!} + \cdots,\ x\mathbb{\in R}\]

Note that \(\sin(x) = x - \frac{x^{3}}{3!} + \frac{x^{5}}{5!}\)

\[\sin(2x) = 2x - \frac{(2x)^{3}}{3!} + \frac{(2x)^{5}}{5!}\ \]

\[e^{- x} = 1 + ( - x) + \frac{( - x)^{2}}{2!} + \frac{( - x)^{3}}{3!} + \cdots = 1 - x + \frac{x^{2}}{2!} - \frac{x^{3}}{3!} + \cdots\]

\begin{align*}
e^{x} - e^{- x} &= 2x + \frac{2x^{3}}{3!} + \cdots \\
\frac{e^{x} - e^{- x}}{2} &= x + \frac{x^{3}}{3!} + \frac{x^{5}}{5!} + \cdots = \sinh(x) \\
e^{x} + e^{- x} &= 2 + 2\left( \frac{x^{2}}{2!} \right) + 2\left( \frac{x^{4}}{4!} \right) + \cdots \\
\frac{e^{x} + e^{- x}}{2} &= 1 + \frac{x^{2}}{2!} + \frac{x^{4}}{4!} = \cosh(x)
\end{align*}



\subsubsection{Euler's formula}

\begin{align*}
e^{ix} &= 1 + ix + \frac{\left( ix \right)^{2}}{2!} + \frac{\left( ix \right)^{3}}{3!} + \cdots \\
&= 1 + ix - \frac{x^{2}}{2!} - \frac{ix^{3}}{3!} + \cdots
\end{align*}


Note that

\[i \cdot \sin(x) = ix - \frac{ix^{3}}{3!} + \frac{ix^{5}}{5!} - \cdots\]

So:

\[\cos(x) + i \cdot \sin(x) = 1 + ix + \frac{\left( ix \right)^{2}}{2!} + \frac{\left( ix \right)^{3}}{3!} + \cdots\]

Note:

\begin{align*}
\frac{\left( e^{ix} + e^{- ix} \right)}{2} &= \cos(x) \\
\frac{e^{ix} - e^{- ix}}{2i} &= \sin(x)
\end{align*}



\subsection{Taylor series expansion}

Find the Taylor series centered at \(a = 0\):

\[\frac{1}{x^{2} - 3x + 2}\]

We'll have to do partial fractions. If the quadratic is irreducible,
then you can complete the square (don't try otherwise).
\(1 + x^{2} = 0\) is the quintessential irreducible quadratic, where its
roots are \(\pm i\).

\begin{align*}
\frac{1}{x^{2} - 3x + 2} &= \frac{1}{(x - 1)(x - 2)} = \frac{A}{x - 1} + \frac{B}{x - 2} \\
&= \frac{A}{x - 1} \cdot \frac{x - 2}{x - 2} + \frac{B}{x - 2} \cdot \frac{x - 1}{x - 1} \\
0x + 1 &= (A + B)x + ( - 2A - B) \\
A + B &= 0,\  - 2A - B = 1 \\
A &= - B \\
&\vdots  \\
B &= 1,\ A = - 1 \\
&- \frac{1}{x - 1} + \frac{1}{x - 2} \\
&= \frac{1}{1 - x} - \frac{1}{2 - x} = \frac{1}{1 - x} - \frac{1}{2} \cdot \frac{1}{1 - \frac{x}{2}} \\
&= \sum_{n = 0}^{\infty}x^{n} - \frac{1}{2}\sum_{n = 0}^{\infty}\left( \frac{x}{2} \right)^{n}\ \text{given}\ |x| < 1\ \text{and\ }\left| \frac{x}{2} \right| < 1 \Rightarrow |x| < 1
\end{align*}


The series will be the following above as long as the above conditions
are met. If \(x\) is multiplied by anything in \(\frac{1}{1 - x}\), then
\(x\) is swapped with what \(x\) is multiplied by.

\[\frac{\sin(x)}{x}\]

We know:

\begin{align*}
\sin(x) &= x - \frac{x^{3}}{3!} + \frac{x^{5}}{5!} - \cdots \\
\frac{\sin(x)}{x} &= 1 - \frac{x^{2}}{2!} + \frac{x^{4}}{5!} - \cdots
\end{align*}


Now, for
\(\ln(1 - x) = - \sum_{n = 1}^{\infty}\frac{x^{n}}{n},\ |x| < 1\)

So what is \(\ln(1 + x)\)

\[\sum_{n = 1}^{\infty}\frac{( - 1)^{n + 1}x^{n}}{n}\]

So what about:

\[f(x) = \int_{0}^{x}{\ e^{- t^{2}}\text{d}t}\]

So \(f'(x) = e^{- x^{2}}\). If we look for the Taylor series for
\(e^{- x^{2}}\) and integrate it, we will get back \(f\).

\begin{align*}
e^{- x^{2}} &= \sum_{n = 0}^{\infty}\frac{x^{n}}{n!} \\
&\int_{0}^{x}{\ e^{t^{2}}\text{d}t}  \\
&= \int_{0}^{x}{\sum_{n = 0}^{\infty}\frac{t^{n}}{n!}}\text{d}t = f(x)  \\
&= C + \sum_{n = 0}^{\infty}{\int_{0}^{x}{\frac{( - 1)^{n}t^{2n}}{n!}\text{d}t}} \\
&= C + \sum_{n = 0}^{\infty}\frac{( - 1)^{n}x^{2n + 1}}{(n!)(2n + 1)}
\end{align*}


To find \(C\), plug in \(x = 0\).

\begin{align*}
f(0) &= 0 = C + \sum_{n = 0}^{\infty}\frac{( - 1)^{n}0^{2n + 1}}{(n!)(2n + 1)} \\
&= C + \sum_{n = 0}^{\infty}\frac{0}{(n!)(2n + 1)} = C \Rightarrow C = 0
\end{align*}


As the integral from 0 to 0 is 0, always

Meaning

\[\int_{0}^{x}e^{{- t}^{2}} = \sum_{n = 0}^{\infty}\frac{( - 1)^{n}x^{2n + 1}}{(n!)(2n + 1)}\]


\subsection{Convergence or divergence on the boundary of the
interval}

\textbf{The series may converge or diverge on the boundary of the
interval of convergence.} How would we determine if the series converges
or diverges on the boundary?

\textbf{Test for convergence individually on the boundaries, from the
series itself. Do not use ratio tests.}

The moment when the ratio test fails, other steps will have to be taken.
Working with a sum involving factorial, \(p\)-series will not work.

To check whether it converges or diverges is very complicated. For
example:

\[\sum_{n = 3}^{\infty}{( - 1)^{n}(x - 3)^{2n}}\]

Ratio test:

\begin{align*}
\lim_{n \rightarrow \infty}\left| \frac{a_{n + 1}}{a_{n}} \right| &= \lim_{n \rightarrow \infty}\left| \frac{( - 1)^{n + 1}(x - 3)^{2n + 2}}{( - 1)^{n}(x - 3)^{2n}} \right| \\
&= (x - 3)^{2} < 1 \Rightarrow |x - 3| < 1 \Rightarrow 2 < x < 4
\end{align*}


What will always happen, is for \(x\) values in the boundary of the
interval, the series may converge or diverge. That is, when \(x = 2\) or
\(x = 4\). To check for this, we would need to do something else.
Considering \(x = 4\), the ratio will exactly be 1, which may not be
possible to run the ratio test with. First check \(x = 4\). Our series
is of the form:

\[\sum_{n = 3}^{\infty}{( - 1)^{n}(4 - 3)^{2n}} = \sum_{n = 3}^{\infty}{( - 1)^{n}(1)^{2n}}\ \text{diverges}\]

If we keep going to infinity, our total is not going closer to a
particular number. The more we add, the closer and closer we get to 2.
As we continue to add 1 and subtract 1, we are not getting closer to any
number. We are bounding between 0 and 1. This is why it counts as
divergent.

We may plug in \(x = 2\):

\[\sum_{n = 3}^{\infty}{( - 1)^{n}( - 1)^{2n}} = \sum_{n = 3}^{\infty}( - 1)^{n}\ \text{diverges}\]

This means that \((2,\ 4)\) is the interval of convergence.

For this course, boundaries should be inspected, but it may not be done
for future courses.

Another one:

\begin{align*}
&\sum_{n = 1}^{\infty}{n(x - 2)^{n}} \\
\lim_{n \rightarrow \infty}\left| \frac{(n + 1)(x - 2)^{n + 1}}{n(x - 2)^{n}} \right| &= |x - 2| < 1 \\
\Rightarrow |x - 2| &< 1
\end{align*}


Do some tests:

\begin{align*}
x &= 1 \\
\sum_{n = 1}^{\infty}{n( - 1)^{n}},\lim_{n \rightarrow \infty}n &\neq 0 \Rightarrow \text{diverges}
\end{align*}


\[x = 3 \Rightarrow \sum_{n = 1}^{\infty}n \Rightarrow \text{diverges}\]

Let's change the question up:

\[\sum_{n = 1}^{\infty}{\frac{1}{n}(x - 2)^{n}\ }\]

What would change? Not much.

\begin{align*}
\lim_{n \rightarrow \infty}\left| \frac{(x - 2)^{n + 1}n}{(n + 1)(x - 2)^{n}} \right| &= |x - 2| < 1 \\
1 &< x < 3
\end{align*}


Check \(x = 3\)

\[\sum_{n = 1}^{\infty}{\frac{1}{n}(1)^{n}\ } \Rightarrow \text{diverges}\]

Check \(x = 1\)

\begin{align*}
&\sum_{n = 1}^{\infty}\left| \frac{1}{n}( - 1)^{n} \right| = \sum_{n = 1}^{\infty}\frac{1}{n} \Rightarrow \text{div.} \\
&\sum_{n = 1}^{\infty}{\frac{1}{n}( - 1)^{n}} \Rightarrow \text{cond.}\ \text{conv.\ }
\end{align*}


The interval of convergence is \(\lbrack 1,\ 3)\)

The same question, but:

\begin{align*}
&\sum_{n = 1}^{\infty}{\frac{1}{n^{2}}(x - 2)^{n}} \\
&\lim_{n \rightarrow \infty}\left| \frac{1}{(n + 1)^{2}}(x - 2)^{n + 1} \cdot \frac{n^{2}}{(x - 2)^{n}} \right| \\
\lim_{n \rightarrow \infty}|x - 2| &= |x - 2| < 1 \\
1 &< x < 3
\end{align*}


Testing boundaries:

\begin{align*}
x &= 3 \\
&\sum_{n = 1}^{\infty}{\frac{1}{n^{2}}(3 - 2)^{n}} \Rightarrow \text{conv} \\
&\sum_{n = 1}^{\infty}{\frac{1}{n^{2}}( - 1)^{n}}
\end{align*}


\[\sum_{n = 1}^{\infty}\frac{1}{n^{2}} \Rightarrow \text{conv}\]

If a series is absolutely convergent, it is convergent. The interval of
convergence is \(\lbrack 1,\ 3\rbrack\).

Anything out of the interval of convergence, the sum doesn't exist, so
there's no approximation going on.


\subsection{Interval of convergence for the geometric
series}

If we try to plot \(\frac{1}{1 - x}\), we have a function that looks
like


\begin{figure}[H]
\centering
\includegraphics[width=1.45607in,height=1.10493in]{latex_images_Sequences_20/media/image4.png}
\end{figure}


We can only approximate the function up to \(|x| < 1\). The function
approaching infinity gives us the radius of convergence. In the real
plane and the complex plane, anything inside the circle formed by the
radius of convergence will converge.


\section{Taylor series at a given
value}

To find the Taylor series of the functions centered at the given value
of \(a\), you'll need \(x - a\).

\begin{align*}
e^{x} &= e^{(x - a) + a} = e^{x - a} \cdot e^{a} \\
\Rightarrow e^{x} &= \sum_{n = 0}^{\infty}\frac{x^{n}}{n!} = \sum_{n = 1}^{\infty}{\frac{(x - a)^{n}}{n!} \cdot e^{a}} \\
&= \sum_{n = 0}^{\infty}{\frac{e^{a}}{n!}(x - a)^{n}}
 \\
\frac{1}{1 - x} &= \frac{1}{1 - \left( (x - 3) + 3 \right)} = \frac{1}{- 2} \cdot \frac{1}{1 - \frac{x - 3}{- 2}} \\
&= - \frac{1}{2}\sum_{n = 0}^{\infty}\left( \frac{x - 3}{- 2} \right)^{n}
\end{align*}


Also: \(\left| \frac{x - 3}{- 2} \right| < 1 \Rightarrow |x - 3| < 2\)

Note that \(x\) will work for every interval as exponentials'
intervals\ldots{}


\subsection{To deal with a constant}

\[f(x) = f\left( (x - a) + a \right)\]

Example with \(\sin(x)\) at \(a = \frac{\pi}{4}\). The sum of angles
formula must be used:

\begin{align*}
\sin(x) &= \sin\left( \left( x - \frac{\pi}{4} \right) + \frac{\pi}{4} \right) \\
&= \sin\left( x - \frac{\pi}{4} \right)\cos\left( \frac{\pi}{4} \right) + \cos\left( x - \frac{\pi}{4} \right)\sin\left( \frac{\pi}{4} \right) \\
&= \sum_{n = 0}^{\infty}{\frac{( - 1)^{n}}{(2n + 1)!}\left( x - \frac{\pi}{4} \right)^{2n + 1} \cdot \frac{1}{\sqrt{2}}} + \sum_{n = 0}^{\infty}{\frac{- 1}{(2n)!}\left( x - \frac{\pi}{4} \right)^{2n}} \cdot \frac{1}{\sqrt{2}}
\end{align*}



\subsection{Exact value of some
series}

What does \(\sum_{n = 1}^{\infty}\frac{1}{n^{2}}\) converge to?
\(\frac{\pi^{2}}{6}\). We'll never get this result.

The geometric series formula is:

\[\frac{1}{1 - x} = \sum_{n = 0}^{\infty}x^{n}\]

Which happens to be the Taylor series. If you encounter a series in the
form

\[\sum_{n = 1}^{\infty}{nx^{n}}\]

Work backwards.

\[\frac{1}{1 - x} = \sum_{n = 0}^{\infty}x^{n}\]

Take the derivative of both sides:

\[- \frac{1}{(1 - x)^{2}} \cdot ( - 1) = \sum_{n = 1}^{\infty}{nx^{n - 1}}\]

Keep in mind:
\(\frac{d}{\text{d}x}\sum_{n = 0}^{\infty}x^{n} = \frac{d}{\text{d}x}\left( 1 + x + x^{2} + x^{3} \right) = 1 + 2x + 3x^{2} + 4x^{3} + \cdots = \sum_{n = 1}^{\infty}{nx^{n - 1}}\)
meaning we can get rid of the starting term. That is when the former
starting term ends up as 0 after the derivative is taken.

Multiply \(x\) on both sides:

\begin{align*}
\frac{x}{(1 - x)^{2}} &= x\sum_{n = 1}^{\infty}{nx^{n - 1}} \\
\frac{x}{(1 - x)^{2}} &= \sum_{n = 1}^{\infty}{nx^{n}}
\end{align*}


There we go.

Another one:

\[\sum_{n = 0}^{\infty}\frac{x^{n + 1}}{n!}\]

\begin{align*}
e^{x} &= \sum_{n = 0}^{\infty}\frac{x^{n}}{n!} \\
xe^{x} &= \sum_{n = 0}^{\infty}\frac{x^{n + 1}}{n!!}
\end{align*}


Another one:

\begin{align*}
&\sum_{n = 0}^{\infty}\frac{x^{n}}{(n + 2)n!} \\
e^{x} &= \sum_{n = 0}^{\infty}\frac{x^{n}}{n!} \\
xe^{x} &= \sum_{n = 0}^{\infty}\frac{x^{n + 1}}{n!}
\end{align*}


Take the derivative of both sides with respect to \(x\):

\[e^{x} + xe^{x} = \sum_{n = 0}^{\infty}\frac{(n + 1)x^{n}}{n!}\]

Another example to get \(n + \text{something}\) on the denominator:

\[\sum_{n = 0}^{\infty}\frac{x^{n}}{(n + 2)n!}\]

\begin{align*}
e^{x} &= \sum_{n = 0}^{\infty}\frac{x^{n}}{n!} \\
xe^{x} &= \sum_{n = 0}^{\infty}\frac{x^{n + 1}}{n!}
\end{align*}


\[xe^{x} - e^{x} = C + \sum_{n = 0}^{\infty}\frac{x^{n + 2}}{(n + 2)n!}\]

When you integrate both sides, add the constant of integration. Most of
the time, it is 0, but best to check it every time by setting \(x\) to
0.

\(x = 0\):

\begin{align*}
0e^{0} - e^{0} &= C + \sum_{n = 0}^{\infty}{\frac{1}{n + 2} \cdot \frac{0^{n + 2}}{n!}} = C + 1 + 0 \\
- 1 &= C
\end{align*}


Be a bit careful when you think sums evaluate to zero -- there may be
causes where something might be raised to the power of zero. You may
need to expand sums first.

Some things:

\begin{align*}
xe^{x} - e^{x} + 1 &= \sum_{n = 0}^{\infty}\frac{x^{n + 2}}{(n + 2)n!} \\
\frac{xe^{x} - e^{x} + 1}{x^{2}} &= \sum_{n = 0}^{\infty}\frac{x^{n}}{(n + 2)n!}
\end{align*}


Hardest example:

\begin{align*}
\sin(x) &= \sum_{n = 0}^{\infty}{\frac{( - 1)^{n}}{(2n + 1)!}x^{2n + 1}} \\
x\sin(x) &= \sum_{n = 0}^{\infty}{\frac{( - 1)^{n}}{(2n + 1)!}x^{2n + 2}} \\
\sin(x) + x\cos(x) &= \sum_{n = 0}^{\infty}{\frac{( - 1)^{n}}{(2n + 1)!}(2n + 2)x^{2n + 1}}
\end{align*}



\subsection{Finding limits using Taylor
series}

\begin{align*}
&\lim_{x \rightarrow 0}\frac{6\sin(x) - 6x + x^{3}}{x^{5}} \\
&= \lim_{x \rightarrow 0}\frac{6\left( x - \frac{x^{3}}{3!} + \frac{x^{5}}{5!} - \frac{x^{7}}{7!}\cdots \right) - 6x + x^{3}}{x^{5}} \\
&= \lim_{x \rightarrow 0}\frac{6\left( \frac{x^{5}}{5!} - \frac{x^{7}}{7!}\cdots \right)}{x^{5}} \\
&= \lim_{x \rightarrow 0}{\frac{6}{5!} - \frac{6x^{2}}{7!} + \frac{6x^{4}}{9!} - \cdots} \\
&= \frac{6}{5!} = \frac{1}{20}
\end{align*}


Because \(x = 0\), all terms that are multiplied by \(x\) cancel out.

Another:

\begin{align*}
&\lim_{x \rightarrow 0}\frac{e^{x^{2}} - \cos(2x) - 3x^{2}}{x^{2}\sin\left( x^{2} \right)} \\
&= \lim_{x \rightarrow 0}\frac{\left( 1 + x^{2} + \frac{x^{4}}{2} + \frac{x^{6}}{3!} + \cdots \right) - \left( 1 - \frac{(2x)^{2}}{2!} + \frac{(2x)^{4}}{4!} - \cdots \right) - 3x^{2}}{x^{2}\left( x^{2} - \frac{x^{6}}{3!} + \frac{x^{10}}{5!} - \cdots \right)}
\end{align*}



\section{Telescoping series}

A \textbf{telescoping series} is when adjacent positive and negative
terms cancel so the \(N\)th partial sum: \(\sum_{n = 1}^{N}a_{n}\) would
only consist of the first and the last term, and we can compute the
series by taking the limit \(N \rightarrow \infty\).

\[\sum_{n = 1}^{\infty}a_{n} = \lim_{N \rightarrow \infty}{\sum_{n = 1}^{N}a_{n}}\]

Example:

\begin{align*}
&\sum_{n = 1}^{\infty}{\arctan(n) - \arctan(n + 1)} = \lim_{N \rightarrow \infty}{\sum_{n = 1}^{N}{\arctan(n) - \arctan(n + 1)}} \\
&= \arctan(0) - \arctan(1) + \arctan(1) - \arctan(2) + \arctan(2) - \arctan(3) + \cdots + \arctan(N) - \arctan(N + 1)
\end{align*}


Because we're adding one positive and one negative, the adjacent terms
cancel out. We're left out with the first and last term:

\begin{align*}
&= \arctan(0) - \arctan(N + 1) \\
&= - \frac{\pi}{2}
\end{align*}



\section{Power series}

\[g(x) = \sum_{n = 0}^{\infty}{a_{n}(x - a)^{n}}\]

is called a power series.

For any power series \(g(x)\), there exists a unique radius of
convergence \(R \geq 0\).

\begin{enumerate}
\def\labelenumi{\arabic{enumi}.}
\item
  The power series \(g(x)\) converges absolutely for \(|x - a| < R\) (in
  the circle).
\item
  The power series \(g(x)\) diverges for \(|x - a| > R\) (outside the
  circle).
\item
  You must check whether the series diverges or converges on the
  boundaries.
\end{enumerate}

If \(R > 0\), then the power series \(g(x)\), as a function of \(x\), is
analytic in \(|x - a| < R\) and \(g'(x)\) is given by term-by-term
differentiation:

\[g'(x) = \sum_{n = 1}^{\infty}{na_{n}(x - a)^{n - 1}}\]

Within the same radius of convergence \(R\).


\subsection{Solving for the radius of
convergence}

\[R = \frac{1}{\lim_{n \rightarrow \infty}\left| \frac{a_{n + 1}}{a_{n}} \right|}\]


\subsection{Taylor polynomial}

Why is the series equal to the function? If we only take finitely many
terms of the Taylor Series, we get a polynomial called the Nth order
polynomial. The Taylor polynomial will not be equal to \(f(x)\), but it
will be an approximation. The difference between the Taylor Polynomial
and \(f(x)\) which is called the remainder \(R_{N}(x)\).

\begin{align*}
P_{N}(x) &= \sum_{n = 0}^{N}{\frac{f^{(n)}(a)}{n!}(x - a)^{n}} \\
f(x) &= \sum_{n = n}^{N}{\frac{f^{(n)}(a)}{n!}(x - a)^{n}} + R_{N}(x) \\
R_{N}(x) &= \frac{f^{(N + 1)}(c)}{(N + 1)!}(x - a)^{N + 1}
\end{align*}


Why does adding up a lot of terms give you a precise quantity? Because
adding more terms give you a more precise approximation. Because \(n!\)
grows extremely quickly, as long as your \(N + 1\)th derivative doesn't
get too big, it will grow slower than \((N + 1)!\) -- since \((N + 1)!\)
is extremely large, all other values:
\(f^{(N + 1)}(c) \cdot (x - a)^{N + 1}\) will be small.


\subsection{First order Taylor
polynomial}

The first-order Taylor polynomial is the tangent line at \(a\).

\[P_{1}(x) = f(a) + f'(a)(x - a)\]

Why is this first-order Taylor polynomial giving the best approximation?
Why couldn't we give the line a different slope?

The definition of the derivative,
\(f'(a) = \lim_{x \rightarrow a}\frac{f(a) - f(a)}{x - a}\) gives:

\begin{align*}
\lim_{x \rightarrow a}\frac{f(x) - P_{1}(x)}{x - a} &= \lim_{x \rightarrow a}\frac{f(x) - f(a) - f'(a)(x - a)}{x - a} \\
&= f'(a) - f'(a) = 0
\end{align*}


\textbf{This is why the tangent line is good at making approximations.}

Let's say I want to use the second-order Taylor approximation. Why is
the second order Taylor approximation more special than any other
quadratic? We can see that the error we're making is tending towards 0?

\[\lim_{x \rightarrow a}\frac{f(x) - P_{2}(x)}{(x - a)^{2}} = 0\]

The tangent line is a good approximation of \(y = f(x)\) near
\(\left( a,\ f(a) \right)\):

\[y - f(a) = f'(a)(x - a)\]

As \(N\) increases, the speed for which \((x - a)^{N}\) is going to zero
increases. In general, the Nth order Taylor Polynomial would have its
error go to zero \textbf{faster than N powers} \((x - a):\)

\[\lim_{x \rightarrow a}\frac{f(x) - P_{N}(x)}{(x - a)^{N}} = 0\]


\begin{figure}[H]
\centering
\includegraphics[width=1in,height=0.8642in]{latex_images_Sequences_20/media/image5.png}
\end{figure}

\end{document}
